\documentclass[a4paper]{amsart}

\usepackage{color}
\usepackage[colorlinks=true, citecolor=red]{hyperref}  %needs to come before amsrefs package to work with \cite{}*{}
\usepackage{amssymb,latexsym}
\usepackage[alphabetic]{amsrefs}
\usepackage[mathscr]{eucal}
\usepackage{pdfsync}
\usepackage{graphicx}
\usepackage{epstopdf}
\DeclareGraphicsRule{.tif}{png}{.png}{`convert #1 `dirname #1`/`basename #1 .tif`.png}
\usepackage[all]{xy}
\usepackage{titletoc}

\CompileMatrices

\theoremstyle{plain}
\newtheorem{theorem}{Theorem}[section]
\newtheorem{lemma}[theorem]{Lemma}
\newtheorem{proposition}[theorem]{Proposition}
\newtheorem{corollary}[theorem]{Corollary}

\theoremstyle{definition}
\newtheorem{definition}[theorem]{Definition}%[section]

\theoremstyle{remark}
\newtheorem{example}[theorem]{Example}
\newtheorem{remark}[theorem]{Remark}
\newtheorem{discussion}[theorem]{Discussion}
\DeclareMathOperator{\Lip}{Lip}
\DeclareMathOperator{\dist}{dist}
\DeclareMathOperator{\expected}{\mathbb{E}}
\DeclareMathOperator{\prb}{Prob}
\DeclareMathOperator{\spt}{spt}
\DeclareMathOperator{\lap}{\Delta}

\raggedbottom

\addtolength{\textwidth}{2cm}
\addtolength{\oddsidemargin}{-1cm}
\addtolength{\evensidemargin}{-1cm} 
\addtolength{\topmargin}{-1cm} 
\addtolength{\textheight}{2cm}
% \renewcommand{\l@subsection}{\@contentsline{2}{5em}{7em}}

%\makeatletter
%\makeatother

 
%%% BEGIN DOCUMENT
\begin{document}

\title{Natural classes of matrices in Euclidean space}  
\author{John E.\ Hutchinson}
\date{\today}

\maketitle


%\renewcommand*\l@subsection{\@tocline{4}{50em}{7em}{10em}{10em}}


\tableofcontents

I will summarise the main results for linear operators $ T: \mathbb{C}^n \to \mathbb{C}^n $ 
( $ T: \mathbb{R}^n \to \mathbb{R}^n $) which are related to the inner product on $ \mathbb{C}^n$ ($ \mathbb{R}^n$). 

Often I will confuse the operator and the matrix of the operator with respect 
to some basis (the basis being understood from context). 

If the proofs are not relatively straightforward I will indicate this. 
%%%%%%%%%%%%%%%%%%%%%%%%%%%%%%%%%%%%%%
%%%%%%%%%%%%%%%%%%%%%%%%%%%%%%%%%%%%%%%%
%%%%%%%%%%%%%%%%%%%%%%%%%%%%%%%%%%%%%%%%%%
\section{\textbf{Existence of a Good Orthonormal Basis}}

\subsection{Self-adjoint case}
The operator  $ T $  is \emph{self-adjoint/Hermitian} iff it corresponds to a (real) stretch in $ n $ 
orthogonal directions ($ \mathbb{C} $  and $ \mathbb{R} $  case). 
Note that in the $ \mathbb{C} $  case, a ``direction'' corresponds to a $ 1 $-complex-dimensional 
subspace and hence to a $ 2  $-real-dimensional subspace. A stretch corresponds to multiplying by a real number in each of these $ n $  complex subspaces. 

Equivalently, $ T $  is self-adjoint iff there is an orthonormal basis w.r.t.\ which 
$ T $  is a real diagonal matrix. 
\[
T = 
   \begin{bmatrix} % or pmatrix or bmatrix or Bmatrix or ...
      \lambda_1  & & \\
       & \ddots & \\
       & & \lambda_n 
   \end{bmatrix}
   \]
Think of the unit ball at the origin being stretched into an ellipsoid. 


\subsection{Unitary case}
$ U $  is \emph{unitary} ($ \mathbb{C} $  case) iff it corresponds to a rotation (multiplication by 
$ e^{i\theta_1}, \dots, e^{i\theta_1} $) in $ n $  orthogonal directions $ v_1,\dots, v_n $.
Equivalently, $ U $  is unitary iff there is an orthonormal basis w.r.t.\ which $ U $  
is a diagonal matrix with all diagonal entries having absolute value 1. 
\[
U = 
   \begin{bmatrix} % or pmatrix or bmatrix or Bmatrix or ...
      e^{i\theta_1}  & & \\
       & \ddots & \\
       & & e^{i\theta_n} 
   \end{bmatrix}
   \]

The corresponding result for orthonormal matrices ($ \mathbb{R} $  case) is not so clean 
and it is better to define orthonormal matrices as we do in ***. The result 
is that $ O $  is orthonormal iff there is an orthonormal basis with respect to 
which $ O $ has matrix 
\[
\setcounter{MaxMatrixCols}{12} 
O = 
   \begin{bmatrix}  % or pmatrix or bmatrix or Bmatrix or ...
      1 & & & & & & & && &&  \\
       & \ddots &&&&&&&&&& \\
       & & 1 &&&&&&&&& \\
       &&& -1 &&&&&&& \\
       &&&& \ddots &&&&&&\\
       &&&& &-1 &&&&&&\\
       &&&&&& \cos \theta_1 & - \sin\theta_1 &&&&\\
       &&&&&& \sin\theta_1 & \cos \theta_1 &&&&\\
       &&&&&&&& \ddots &&&\\
       &&&&&&&&& \cos \theta_k & - \sin\theta_k&\\
       &&&&&&&&& \sin\theta_k & \cos \theta_k
   \end{bmatrix}
   \]
Thus there is an orthonormal basis such that $ O $ fixes some vectors, reflects 
others, and rotates in orthogonal planes determined by pairs of the 
remaining basis vectors. In particular, in $ \mathbb{R}^3 $  every orthonormal operator 
is a rotation around some vector followed perhaps by a reflection in the 
plane orthogonal to this vector! 

Think of the unit ball at the origin being rotated into itself (perhaps there 
is also a reflection, thus reversing the orientation)


\subsection{Normal case}
$ T $  (in the $ \mathbb{C} $  case) is normal iff it corresponds to multiplication by a  complex numbers in each of $ n $  orthogonal directions. That 
is, to a stretch followed by a rotation in each of the corresponding $ n $  2-real dimensional subspaces. 
Equivalently, iff  there is an orthonormal basis with respect to which $ T $  is 
diagonal (with complex, not necessarily real, entries). 

%%%%%%%%%%%%%%%%%%%%%%%%%%%%%%%%%%%%%%%%%%%
\section{\textbf{Adjoints}}
The adjoint $ T^* $  of the operator $  T $ is defined by $ \langle T^* u,v \rangle  = \langle u, T v\rangle $  
for all $ u, v \in \mathbb{C} ^n\  (\mathbb{R}^n )$.

With respect to any orthonormal basis, the matrix $ T^* $ is the conjugate (not 
necessary of course in the $ \mathbb{R} $  case) of the transpose of $ T $, i.e.\  $ T^*_{ij} = \overline{T_{ji}}$.
 

%%%%%%%%%%%%%%%%%%%%%%%%%%%%%%%%%%%%%%%%%%%
\section{\textbf{Analytic and Matrix Approaches}}

\subsection{Self-adjoint case}
$ T$  is self adjoint iff $ T = T^* $. 

Because of what this means for matrices, one often says $ T$  is Hermitian 
in the $ \mathbb{C}$   case and symmetric in the $ \mathbb{R} $  case. 

The  $ \Longrightarrow $ case is easy but the converse requires work. It follows easily from 
the more general result for normal operators below 


\subsection{Unitary case} $ U $ is unitary ($ O $ is orthonormal) iff $ U U^*  = U^* U = I$  ($ OO^* = O^*O = I $). 
(Clear from the matrix version.) 

For matrices this is equivalent to the columns being orthonormal (norm 
one and inner products of different columns equalling zero) and also equivalent to the rows having the same property. 

\subsection{Normal case} $ T$  is normal iff  $ T T^* = T^* T $. ($ \Longrightarrow$   is clear from the definition. The converse  requires work). 



%%%%%%%%%%%%%%%%%%%%%%%%%%%%
\section{\textbf{Polar Decomposition}}

Every linear operator $ T$  ($ \mathbb{C} $  case) can be written in the form $ T = U H $ 
where $ U$  is unitary and $ H  $ is Hermitian (i.e. self-adjoint) with nonnegative 
e-values. That is, real nonnegative stretching in orthogonal complex directions 
followed by rotations in another set of orthogonal complex directions. 

(If the rotations are in the same complex directions as the stretches, then 
the matrix is clearly normal.) 

\medskip
 In the $ \mathbb{R} $  case $ T = OS$  where $ O$  is orthonormal and $ S$  is symmetric with 
nonnegative e-values. That is, $ T$  is again a nonnegative stretch in orthogonal directions and then apply orthogonal transformation. 

\medskip Think of taking the unit ball sitting at the origin, stretching it to an ellipsoid, 
and then rotating. 



%%%%%%%%%%%%%%%%%%%%%%%%%%%%
\section{\textbf{Another Characterisation}}


 $ T$  is selfadjoint iff $ \langle T u, u\rangle $  is real for all u 

$ T $ is unitary (orthonormal) iff  $ \| T u\| = \|T^* u\| = \|u\| $  for all $ u $ 

$ T $ is normal iff  $ \|T u\| = \|T^* u\| $ for all $ u $.

\medskip The proofs are fairly straightforward (with the HINT from the relevant assignment problem). 


\end{document}




%%%%%%%%%%%%%%%%%%%%%%%%%%%%%%%%%%%%%%%%%%%%%%%
%%%%%%%%%%%%%%%%%%%%%%%%%%%%%%%%%%%%%%%%%%%%%%%%%%
%%%%%%%%%%%%%%%%%%%%%%%%%%%%%%%%%%%%%%%%%%%%%%%

\begin{thebibliography}{9} 



\bib{RN}{book}{
   author={Riesz, Frigyes},
   author={Sz.-Nagy, B{\'e}la},
   title={Functional analysis},
   publisher={Dover Publications Inc.},
   note={Translated from the second French edition by Leo F. Boron;
   Reprint of the 1955 original},
   date={1990},
   pages={xii+504},
}
		



\end{thebibliography}





\end{document}


